\minitoc 

% ========================================================================================================
% Cadre Général

\section{Cadre Général}

Dans ce chapitre, nous allons présenter un algorithme permettant de déterminer les solutions de l'équation :
    $$ h( \theta) = 0 $$ 
pour une fonction $h$ lorsqu'on peut écrire $ h ( \theta) = \mathbb{E}(g( \theta,Y))$. 
Par exemple, chercher $ \theta$ tel que : 
    \begin{align*}
        &\int_{\R_+} (x - \theta) e^{-x^3} \; dx = 0 \\ 
        \iff & \int_{\R_+} x e^{-x^3} \; dx - \theta \int_{\R_+} e^{-x^3} \; dx = 0 \\ 
        \iff & \theta = \frac{ \int_{\R_+} x e^{-x^3} \; dx }{ \int_{\R_+} e^{-x^3} \; dx }
    \end{align*}

% \begin{example}
%     \begin{itemize} 
%         \item si $g( \theta, Y) = \theta - Y$ alors $ h( \theta) = 0 \iff \theta = \mathbb{E}(Y)$.
%         \item $g( \theta, Y) = \alpha - \mathbbm{1}_{ Y \leqslant \theta}$ alors $h( \theta) = 0 $ 
%         ssi $ \theta$ est le quantile d'ordre $ \alpha$ de la loi $Y$.   
%     \end{itemize}
% \end{example}

\begin{proposition}[Algorithme]
    On se fixe une suite $ ( \gamma_n)_{n \in \N} \in (\R_+^*)^\N$ telle que 
    $ \sum_{n \geqslant 0} \gamma_n = \infty$ mais $ \sum_{n \geqslant 0} \gamma_n^2 < \infty$. 
    On définit par récurrence : 
        \begin{equation}
            \boxed{\theta_{n+1} = \theta_n - \gamma_n g ( \theta_n, Y_{n+1})}
        \end{equation}
    où les $(Y_n)_{n \in \N}$ sont des variables aléatoires idd de même loi que $Y$.
    Sous de bonnes conditions, $ (\theta_n)_{n \in \N}$ converge vers une solution $ \theta$ de 
    $ h( \theta) = 0$.
\end{proposition}

\begin{example}
    Supposons que $ g( \theta, Y) = \theta - Y$ et $ \gamma_n = \frac{1}{n+1}$.
    La solution de $h( \theta) = 0$ est exactement $ \theta = \mathbb{E}(Y)$.
    On remarque que : 
        \begin{align*}
            \theta_{n+1} &= \theta_n - \frac{1}{n+1} \left( \theta_n - Y_{n+1} \right) \\ 
            &= \frac{n}{n+1} \theta_n + \frac{1}{n+1} Y_{n+1} 
        \end{align*}
    Donc $ (n+1) \theta_{n+1} - n \theta_n = Y_{n+1}$, d'où $ n \theta_n = Y_1 + \dots + Y_n$.
    Soit $ \theta_n = \frac{1}{n } \sum_{i=1}^{n} Y_i$. On a bien $ \lim_{n \to \infty} \theta_n = \mathbb{E}(Y)$ 
    d'après la loi des grands nombres.
\end{example}

% \subsection{Cas d'une fonction déterministe $h$}

% On suppose ici que l'on connaît la fonction $ h : \R \longrightarrow \R$ que l'on 
% sait calculer. On souhaite aussi que $h \in \mathcal{C}^1(\R)$.

% \begin{proposition}[Dichotomie]
%     Pour trouver les solutions de l'équation $h ( \theta) = 0$, on peut essayer une 
%     méthode par dichotomie. On calcule plusieurs valeurs de $h$ jusqu'à trouver une valeur 
%     positive et une négative. Par continuité, on sait que la fonction s'annule entre les deux. 
%     Or cette méthode fonctionne que sur $\R$ et nous devons définir des points de départ 
%     pour l'algorithme.
% \end{proposition}

% \begin{proposition}
%     Nous pouvons aussi considérer une méthod de Newton de la forme : 
%         \begin{equation}
%             \forall n \in \N, \quad x_{n+1} = x_n - \frac{g(x_n)}{g'(x_n)}
%         \end{equation}
%     Nous devons, ici aussi, être capable de très bien calculer $g$ pour des valeurs très 
%     proches pour estimer des valeurs de $g'$. On peut donc proposer une autre méthode : 
%         \begin{equation}
%             \forall n \in \N, \quad x_{n+1} = x_n - \alpha_n g(x_n) 
%         \end{equation}
%     où $ \sum_{i \geqslant 1} \alpha_n = \infty$ et $ \sum_{i \geqslant 1} \alpha_x^2 < \infty$.
%     Alors, la suite $(x_n)_{n \in \N}$ converge bien vers une solution de $g( \theta) = 0$.
% \end{proposition}

% \subsection{Cas d'une fonction non déterministe (cas stochastique)}

% On a donc l'équation de la forme : 
%     \begin{equation}\label{eq:equation-espe}
%         h( \theta) = \mathbb{E}(g( \theta, Y)) = 0
%     \end{equation} 
% On pose donc l'itéré suivant : 
%     $$ \theta_{n+1} = \theta_n - \alpha_n H( \theta_n, Y_{n+1}) $$
% Alors, par construction, $( \theta_n)_{n \in \N}$ est une chaîne de Markov inhomogène 
% dans le temps. On va donc exhiber des conditions sur $H$ pour que $( \theta_n)_{n \in \N}$ 
% converge p.s vers une solution $ \theta^*$ de \ref{eq:equation-espe}.

% \medskip 

% On note donc $ \mathcal{F}_n = \sigma( \theta_0, Y_1, \dots, Y_n)$ la filtration adaptée à 
% $( \theta_n)_{n \in \N}$. On a donc : 
%     \begin{align*}
%         \mathbb{E}( \theta_{n+1} \mid \mathcal{F}_n) &= \mathbb{E}( \theta_n - \alpha_n H( \theta_n, Y_{n+1}) \mid \mathcal{F}_n) \\ 
%         &= \theta_n - \alpha_n \mathbb{E}(H( \theta_n, Y_{n+1}) \mid \mathcal{F}_n) \\ 
%         &= \theta_a - \alpha_n h( \theta_n)
%     \end{align*}

% Donc, conditionnellement à $ \mathcal{F}_n$, $ \mathbb{E}( \theta_{n+1})$ est exactement 
% l'itéré définit par la méthode de Newton. C'est le même principe que l'algorithme de Robbins 
% Monro. On a donc : 
%     \begin{equation}
%         \theta_{n+1} = \theta_n - \alpha_n h( \theta_n) + \alpha_n \Delta M_n 
%     \end{equation}
% où $ \Delta M_n = h( \theta_n) - H( \theta_n, Y_{n+1})$ est une variable centrée 
% indépendante de $\mathcal{F}_n$.

Dans le cas déterministe, on suppose ggg connue et calculable, ce qui permet d’utiliser des méthodes comme Newton. 
Cependant, en pratique, ggg est souvent inconnue (définie par une espérance) ou bruitée (observations aléatoires). 
L’algorithme de Robbins-Monro résout ce problème en remplaçant $ g( \theta) $ par une estimation 
$ H(\theta, Y)$, tout en garantissant la convergence vers la solution de $h(\theta) = 0$ sous des conditions
sur les pas $(\alpha_n)$.

% ========================================================================================================
% Convergence de l'algorithme

\section{Convergence de l'algorithme}

Toute cette section est consacrée à la démonstration de la convergence de l'algorithme. Pour cela, nous allons introduire 
un théorème très important dans l'étude des algorithmes stochastiques : le théorème de Robbins-Sigmund. 

\subsection{Théorème de Robbins-Sigmund}

Nous allons étudier ce théorème sous la forme suivante : 

\begin{theorem}[Robbins-Sigmund]
    Soient quatre suites $(z_n)_{n \in \N}, ( \alpha_n)_{n \in \N}, ( \beta_n)_{n \in \N} $ et $ ( \gamma_n)_{n \in \N}$
    adaptées pour une filtration $( \mathcal{F}_n)_{n \in \N}$. Supposons qu'elles vérifient : 
        \begin{enumerate}[label=\roman*)]
            \item $ \displaystyle \sum_{n \geqslant 0} \alpha_n = \infty $ p.s
            \item $ \displaystyle \sum_{n \geqslant 0} \beta_n < \infty $ p.s 
            \item $ \displaystyle \sum_{n \geqslant 0} \gamma_n < \infty $ p.s
            \item $ \displaystyle \mathbb{E}(z_{n+1} \mid \mathcal{F}_n) \leqslant (1 - \alpha_n + \beta_n) z_n + \gamma_n$
        \end{enumerate}
    
    \medskip
    
    Alors : 
        \begin{equation}
            \lim_{n \to \infty} z_n = 0 \text{ ps}  
        \end{equation}
\end{theorem}

Dans notre situation, nous pouvons considérer que $ \beta_n = 0$ pour tout $n \in \N$. 
En effet, nous cherchons à modifier l'hypothèse $\mathbb{E}(z_{n+1} \mid \mathcal{F}_n) \leqslant (1 - \alpha_n + \beta_n) z_n + \gamma_n$ 
en : 
    $$ \mathbb{E}(z_{n+1} \mid \mathcal{F}_n) \leqslant (1 - \alpha_n) z_n + \gamma_n $$ 

\begin{lemma}
    Réduisons le théorème au cas $ \beta_n = 0$.
\end{lemma}

\begin{proof}
    Supposons donc que : 
    \begin{enumerate}[label=\roman*]
        \item $ \sum_{n \geqslant 0} \alpha_n = \infty $ p.s
        \item $ \sum_{n \geqslant 0} \beta_n < \infty $ p.s 
        \item $ \sum_{n \geqslant 0} \gamma_n < \infty $ p.s
        \item $ \mathbb{E}(z_{n+1} \mid \mathcal{F}_n) \leqslant (1 - \alpha_n) z_n + \gamma_n$
    \end{enumerate}
    Montrons que si $ \beta_n = 0$, alors $ z_n \underset{n \to \infty}{\longrightarrow} 0$ aussi.
    Pour cela, posons, pour tout $n \in \N$ : 
        $$ Y_n = \frac{z_n}{ \prod_{j=1}^n (1 + \beta_j)} $$ 
    On a les égalités suivantes : 
        \begin{align*}
            \mathbb{E} \left( \frac{z_{n+1}}{ \prod_{j=1}^{n+1} (1 + \beta_j)} \mid \mathcal{F}_n\right) &= \frac{ \mathbb{E}(z_{n+1} \mid \mathcal{F}_n)}{\prod_{j=1}^{n+1} (1 + b_j)} \\ 
            & \leqslant \frac{ (1 - \alpha_n + \beta_n) z_n}{ \prod_{j=1}^{n+1} (1 + \beta_j)} + \frac{ \gamma_n}{ \prod_{j=1}^{n+1} (1 + \beta_j)} \\
            &= \frac{1 - \alpha_n + \beta_n}{1 + \beta_{n+1}} \frac{z_n}{\prod_{j=1}^{n} (1 + \beta_j)} + \frac{ \gamma_n}{ \prod_{j=1}^{n+1} (1 + \beta_j)}
        \end{align*}
    Or, pour $n$ grand, on a : 
        $$ \frac{1 - \alpha_n + \beta_n}{1 + \beta_n} \sim_\infty 1 - \alpha_n $$ 
    D'où : 
        $$ \mathbb{E} \left( Y_{n+1} \mid \mathcal{F}_n\right) \leqslant (1 - \alpha_n') Y_n + \gamma_n' $$ 
    Avec : 
        $$ Y_{n} = \frac{z_n}{ \prod_{j=1}^{n} (1 + \beta_j)}, \quad \alpha_n' = \frac{ \alpha_n}{1 + \beta_n} \quad \text{et} \quad \gamma_n' = \frac{ \gamma_n}{\prod_{j=1}^{n} (1 + \beta_j)} $$
    On souhaite maintenant appliquer le theorème de Robbins-Sigmund à la suite $(Y_n)_{n \geqslant 0}$. Pour cela, vérifions 
    les hypothèses du théorème : 
    \begin{enumerate}[label=\roman*]
        \item Pour $ \alpha_n'$, on a que $ \alpha_n' = \frac{ \alpha_n}{1 + \beta_n} $ donc $ \alpha_n' \sim_\infty \alpha_n$. 
        Par comparaison : 
            $$ \sum_{n \geqslant 0} \alpha_n < \infty \text{ p.s } \quad \Longrightarrow \quad \sum_{n \geqslant 0} \alpha_n' < \infty \text{ p.s } $$ 
        \item Pour $ \gamma_n'$, on a : 
            $$ \sum_{n \geqslant 0} \gamma_n' = \sum_{n \geqslant 0} \frac{ \gamma_n}{ \prod_{j=1}^{n} (1 + \beta_j)} < \sum_{n \geqslant 0} \gamma_n < \infty $$ 
    \end{enumerate}
    On peut donc appliquer le théorème d'où : 
        $$ \lim_{n \to \infty} Y_n = 0 \text{ p.s} $$ 
    De plus, par construction de $Y_n$, on a : 
        $$ \lim_{n \to \infty} z_n = \lim_{n \to \infty} Y_n \underbrace{\prod_{j=1}^{n} (1 + \beta_j)}_{< \infty} \quad
            \implies \quad \lim_{n \to \infty} z_n = 0 $$
    Le théorème est donc valable avec ce nouveau jeu d'hypothèses.\qed
\end{proof}

Nous pouvons désormais passer à la démonstration du théorème de Robbins-Sigmund.

\begin{proof}
    La démonstration de ce résultat est basée sur la construction d'une sur-martingale et de l'utilisation 
    du théorème de convergence des martingales pour conclure. Nous commençons donc avec les hypothèses suivantes : 
    \begin{enumerate}[label=\roman*)]
        \item $ \sum_{n \geqslant 0} \alpha_n = \infty $ p.s
        \item $ \sum_{n \geqslant 0} \beta_n < \infty $ p.s 
        \item $ \sum_{n \geqslant 0} \gamma_n < \infty $ p.s
        \item $ \mathbb{E}(z_{n+1} \mid \mathcal{F}_n) \leqslant (1 - \alpha_n) z_n + \gamma_n$
    \end{enumerate}
    Montrons que $ z_n \underset{n \to \infty}{\longrightarrow} 0$.
    
    \medskip

    \emph{1. Construction de la sur-martingale}. 
    Posons : 
        $$ \forall n \in \N, \quad S_n = z_n + \sum_{k=0}^{n-1} \alpha_k z_k - \sum_{k=0}^{n-1} \gamma_k $$ 
    Montrons que $(S_n)_{n \geqslant 0}$ est une sur-martingale. 
    \begin{align*}
        \mathbb{E}(S_{n+1} \mid \mathcal{F}_n) &= \mathbb{E}(z_{n+1} \mid \mathcal{F}_n) + \sum_{k=0}^{n} \mathbb{E}(\alpha_k z_k \mid \mathcal{F}_n)
            - \sum_{k=0}^{n} \mathbb{E}( \gamma_k \mid \mathcal{F}_n) \\ 
        &= \mathbb{E}(z_{n+1} \mid \mathcal{F}_n) + \sum_{k=0}^{n} \alpha_k z_k - \sum_{k=0}^{n} \gamma_k \\ 
        &\leqslant (1 - \alpha_n) z_n + \gamma_n + \sum_{k=0}^{n} \alpha_k z_k - \sum_{k=0}^{n} \gamma_k \\
        &= \sum_{k=0}^{n-1} \alpha_k z_k - \sum_{k=0}^{n-1} \gamma_k + z_n = S_n 
    \end{align*}
    Par définition, $(S_n)_{n \geqslant 0}$ est bien une sur-martingale pour la filtration $ (\mathcal{F}_n)_{n \geqslant 0}$.

    \medskip

    \emph{2. Convergence de la sur-martingale. } Soit $ \tau_c$ le temps d'arrêt suivant : 
        $$ \tau_c = \inf \left\{ n \in \N \mid \sum_{k=0}^{n} \gamma_k \geqslant  \right\} $$ 
    Soit la martingale arrêtée par $ \tau_C$ $(S_{n \wedge \tau_c})_{n \geqslant 0}$, elle est bornée inférieurement par $-C$ car : 
        $$ S_{n \wedge \tau_C} \geqslant - \sum_{k=0}^{ \tau_C -1} \gamma_k \geqslant - C $$ 
    Par théorème, elle converge donc presque sûrement vers une limite finie $W_C$. 
    Puisque $ \sum_{k \geqslant 1} \gamma_k < \infty$ p.s, il existe $C > \sum_{k \geqslant 1} \gamma_k $ tel que 
    $ \tau_c= \infty$ p.s. Donc sur l'évènnement $\{ \tau_c = \infty\}$, on a $S_n = S_{n \wedge \tau_C}$ pour tout $n$, donc :
        $$ \lim_{n \to \infty} S_n = W_c \text{ p.s} $$ 
    Or $S_n = z_n + \sum_{k=0}^{n-1} \alpha_k z_k - \sum_{k=0}^{n-1} \gamma_k$, et comme $ \sum_{k \geqslant 1} \gamma_k $ converge, 
    la convergence de $S_n$ implique celle de : 
        $$ z_n + \sum_{k=0}^{n-1} \alpha_k z_k $$
    mais $ \sum_{k \geqslant 1} \alpha_k = \infty$ donc nécessairement $ \sum_{k \geqslant1} \alpha_k z_k < \infty$. 
    Finalement : 
        $$ z_n \underset{n \to \infty}{\longrightarrow} 0 \text{ p.s} $$
    \qed 
\end{proof}

\subsection{Convergence de l'algorithme de Robbins-Monro}

% Maintenant que nous avons tous les outils nécessaires, nous pouvons passer à l'étude 
% de la convergence de cet algorithme. 

% Pour rappel, on souhaite déterminer les solutions de l'équation $ h( \theta) = 0$ où $h$ peut s'écrire sous la forme : 
%     $$ h( \theta) = \mathbb{E}[H(Y, \theta)] $$ 
% où $Y$ est une variable aléatoire quelconque que l'on sait simuler. 
% L'algorithme de Robbins-Monro approche la solution $ \theta^*$ par les itérés suivants : 
%     $$ \forall n \geqslant 0, \quad \theta_{n+1} = \theta_n - \gamma_n H(Y_{n+1}, \theta_n) $$ 
% où $(Y_n)_{n \geqslant 0}$ est une suite de variables idd de même loi que $Y$. 
% Le théorème suivant nous assure sa convergence sous de bonnes hypothèses : 

% \begin{theorem}[Conditions suffisantes de convergence]
%     Avec les mêmes notations que précédement : 
%     \begin{enumerate}[label=\roman*)]
%         \item (Coercivité) S'il existe $ \theta^*$ tel que $h( \theta) ( \theta - \theta^*) \geqslant \mu \| \theta - \theta^* \|^2$, 
%         \item S'il existe $A, B > 0$ tels que $ \| h( \theta_n) \| ^2 \leqslant A + B \| \theta_n \|^2$, 
%         \item (Contrôle du bruit) Et que $ \mathbb{E}( \| h( \theta) - H( \theta, Y) \|^2) \leqslant K$.
%     \end{enumerate}
%     Alors $\lim_{n \to \infty} \theta_n = \theta^*$ p.s.
% \end{theorem}

% \begin{proof}
%     La démonstration de cette convergence utilise la construction récursive de l'algorithme ainsi 
%     que le théorème de Robbins-Sigmund. 
%     Soit $( \mathcal{F}_n) = \sigma ( \theta_0, \dots, \theta_n, Y_0, \dots, Y_n)$ la filtration associée 
%     au $n$-ième itéré de l'algorithme. L'idée est de séparer l'incrément de l'algorithme en deux valeurs : 
%         $$ \Delta M_n = \underbrace{h( \theta_n)}_{\text{partie déterministe}} - \underbrace{H( \theta_n, Y_(n+1))}_{\text{partie stochastique}} $$
%     noté $\Delta M_n$ et appelée l'incrément de martingale. On aura donc l'algorithme suivant : 
%         $$ \theta_{n+1} = \theta_n - \gamma_n h( \theta_n) + \gamma_n \Delta M_n  $$
%     On utilisera plus tard le fait que $(M_n)_{n \geqslant 0}$ est une martingale d'espérance nulle : 
%         \begin{align*}
%             \mathbb{E}( \Delta M_n \mid \mathcal{F_n}) &= \mathbb{E}(h( \theta_n) \mid \mathcal{F}_n) - \mathbb{E}(H( \theta_n, Y_{n+1}) \mid \mathcal{F}_n) \\ 
%             &= \mathbb{E}(H( \theta_n, Y_{n+1}) \mid \mathcal{F}_n) - \mathbb{E}(H( \theta_n, Y_{n+1}) \mid \mathcal{F}_n) = 0
%         \end{align*}
%     Pour l'étude de la convergence, nous étudierons les distances quadratiques des itérés : 
%         $$ z_n = \| \theta_n - \theta^* \|^2 $$ 
%     ainsi que l'égalité suivante : 
%         $$ \| a b + c \|^2 = \|a\|^2 + \|b\|^2 + \|c\|^2 + 2 \langle a,b \rangle + 2 \langle a, c \rangle + 2 \langle b , c \rangle $$ 
%     On a donc : 
%         \begin{align*}
%             \mathbb{E}(z_{n+1} \mid \mathcal{F}_n) &= \mathbb{E}( \| \theta_n - \theta^* \|^2 \mid \mathcal{F}_n) \\ 
%             &= \mathbb{E}( \| \theta_n - \gamma_n h( \theta_n) + \gamma_n \Delta M_n - \theta^*\|^2 \mid \mathcal{F}_n) \\ 
%             &= \mathbb{E}( \| \underbrace{\theta_n - \theta^*}_{a} \underbrace{- \gamma_n h( \theta_n)}_{b} + \underbrace{\gamma_n \Delta M_n}_{c} \|^2 \mid \mathcal{F}_n) \\ 
%             &= \| \theta_n - \theta^* \| + \gamma_n^2 \| h( \theta_n) \|^2 + \gamma_n^2 \mathbb{E}(\| \Delta M_n \|^2 \mid \mathcal{F}_n) \\
%             & \quad \quad + 2 \gamma_n ( \theta_n - \theta^*)( \Delta M_n - H( \theta_n, Y_{n+1})) - 2 \gamma_n^2 h( \theta_n) \underbrace{\mathbb{E}( \Delta M_n \mid \mathcal{F}_n)}_{ = 0} \\ 
%             &= z_n + \gamma_n^2 \| h( \theta_n) \|^2 + \gamma_n^2 \mathbb{E}( \| \Delta M_n \|^2 \mid \mathcal{F}_n) - 2 \gamma_n ( \theta_n - \theta^*) h( \theta_n) \\
%             & \leqslant z_n + \gamma_n^2 (1 + B \| \theta_n \|^2) + \gamma_n^2 \mathbb{E}( \| \Delta M_n \|^2 \mid \mathcal{F}_n) - 2 \gamma_n \mu \| \theta_n - \theta^*\|^2 \\ 
%             & \leqslant \| \theta_n - \theta^*\|^2 (1 - \underbrace{ 2 \gamma_n \mu}_{ = \alpha_n'}) + \underbrace{\gamma_n^2 ( K + A + B \| \theta_n \|^2)}_{( = \gamma_n')}
%         \end{align*}
%     Et l'on peut conclure grâce au théorème de Robbins-Sigmund.\qed
% \end{proof}


Maintenant que nous avons tous les outils nécessaires, nous pouvons passer à l'étude
de la convergence de cet algorithme.
Pour rappel, on souhaite déterminer les solutions de l'équation $h(\theta) = 0$ où $h$ peut s'écrire sous la forme :
    \[ h(\theta) = \mathbb{E}[H(Y, \theta)], \]
où $Y$ est une variable aléatoire quelconque que l'on sait simuler.
L'algorithme de Robbins-Monro approche la solution $\theta^*$ par les itérés suivants:
    \[ \forall n \geq 0, \quad \theta_{n+1} = \theta_n - \gamma_n H(Y_{n+1}, \theta_n),\]
où $(Y_n)_{n \geq 0}$ est une suite de variables i.i.d.\ de même loi que $Y$.

Le théorème suivant nous assure sa convergence sous de bonnes hypothèses :

\begin{theorem}[Conditions suffisantes de convergence]
    Avec les mêmes notations que précédemment :
    \begin{enumerate}[label=\roman*)]
        \item (Coercivité) Il existe $\theta^*$ tel que $h(\theta)(\theta - \theta^*) \geq \mu \|\theta - \theta^*\|^2$ pour un certain $\mu > 0$,
        \item (Croissance de $h$) Il existe $A, B > 0$ tels que $\|h(\theta)\|^2 \leq A + B \|\theta\|^2$,
        \item (Contrôle du bruit) $\mathbb{E}(\|h(\theta) - H(\theta, Y)\|^2) \leq K$ pour un certain $K > 0$.
    \end{enumerate}
    Alors, si $\sum_{n \geq 0} \gamma_n = \infty$ et $\sum_{n \geq 0} \gamma_n^2 < \infty$, on a :
        \[ \lim_{n \to \infty} \theta_n = \theta^* \quad \text{p.s.}\]
\end{theorem}

\begin{proof}
    La démonstration de cette convergence utilise la construction récursive de l'algorithme ainsi que le théorème de Robbins-Siegmund.
    Soit $\mathcal{F}_n = \sigma(\theta_0, \dots, \theta_n, Y_0, \dots, Y_n)$ la filtration associée au $n$-ième itéré de l'algorithme.
    L'idée est de séparer l'incrément de l'algorithme en deux parties :
    \[
    \Delta M_n = h(\theta_n) - H(\theta_n, Y_{n+1}),
    \]
    appelée incrément de martingale. On peut donc réécrire l'algorithme comme suit :
    \[
    \theta_{n+1} = \theta_n - \gamma_n h(\theta_n) + \gamma_n \Delta M_n.
    \]

    Montrons que $(M_n)_{n \geq 0}$ est une martingale :
    \[
    \begin{aligned}
        \mathbb{E}(\Delta M_n \mid \mathcal{F}_n) &= \mathbb{E}(h(\theta_n) - H(\theta_n, Y_{n+1}) \mid \mathcal{F}_n) \\
        &= h(\theta_n) - \mathbb{E}(H(\theta_n, Y_{n+1}) \mid \mathcal{F}_n) \\
        &= h(\theta_n) - h(\theta_n) = 0.
    \end{aligned}
    \]

    Pour étudier la convergence, nous considérons la distance quadratique des itérés :
    \[
    z_n = \|\theta_n - \theta^*\|^2.
    \]
    ainsi que l'égalité suivante : 
        $$ \| a + b + c \|^2 = \|a\|^2 + \|b\|^2 + \|c\|^2 + 2 \langle a,b \rangle + 2 \langle a, c \rangle + 2 \langle b , c \rangle. $$
    On a donc : 
        \[
            \begin{aligned}
                \mathbb{E}(z_{n+1} \mid \mathcal{F}_n) &= \mathbb{E}(\|\theta_{n+1} - \theta^*\|^2 \mid \mathcal{F}_n) \\
                &= \mathbb{E}(\|\theta_n - \gamma_n h(\theta_n) + \gamma_n \Delta M_n - \theta^*\|^2 \mid \mathcal{F}_n) \\
                &= \|\theta_n - \theta^*\|^2 + \gamma_n^2 \|h(\theta_n)\|^2 + \gamma_n^2 \mathbb{E}(\|\Delta M_n\|^2 \mid \mathcal{F}_n) \\
                &\quad - 2\gamma_n \langle \theta_n - \theta^*, h(\theta_n) \rangle + 2\gamma_n \langle \theta_n - \theta^*, \Delta M_n \rangle \\
                &\quad - 2\gamma_n^2 \langle h(\theta_n), \Delta M_n \rangle.
            \end{aligned}
        \]

    Comme $\mathbb{E}(\Delta M_n \mid \mathcal{F}_n) = 0$, les termes croisés $\langle \theta_n - \theta^*, \Delta M_n \rangle$ et $\langle h(\theta_n), \Delta M_n \rangle$ disparaissent en espérance conditionnelle. On obtient donc:
    \[
        \mathbb{E}(z_{n+1} \mid \mathcal{F}_n) \leqslant z_n + \gamma_n^2 \|h(\theta_n)\|^2 + \gamma_n^2 \mathbb{E}(\|\Delta M_n\|^2 \mid \mathcal{F}_n) - 2\gamma_n \langle \theta_n - \theta^*, h(\theta_n) \rangle.
    \]

    En utilisant les hypothèses du théorème :
    \[
    \begin{aligned}
        \mathbb{E}(z_{n+1} \mid \mathcal{F}_n) & \leqslant z_n + \gamma_n^2 (A + B \|\theta_n\|^2) + \gamma_n^2 K - 2\gamma_n \mu \|\theta_n - \theta^*\|^2 \\
        & \leqslant z_n (1 - 2\gamma_n \mu) + \gamma_n^2 (A + K + B \|\theta_n\|^2).
    \end{aligned}
    \]

    Comme $\|\theta_n\|^2 \leqslant 2\|\theta_n - \theta^*\|^2 + 2\|\theta^*\|^2$, on peut écrire :
    \[
        \mathbb{E}(z_{n+1} \mid \mathcal{F}_n) \leqslant z_n (1 - 2\gamma_n \mu + 2B\gamma_n^2) + \gamma_n^2 (A + K + 2B\|\theta^*\|^2).
    \]

    En posant $\alpha_n = 2\gamma_n \mu - 2B\gamma_n^2$ et $\beta_n = \gamma_n^2 (A + K + 2B\|\theta^*\|^2)$, on obtient :
    \[
        \mathbb{E}(z_{n+1} \mid \mathcal{F}_n) \leqslant (1 - \alpha_n) z_n + \beta_n.
    \]

    Sous les hypothèses $\sum \gamma_n = \infty$ et $\sum \gamma_n^2 < \infty$, on a $\sum \alpha_n = \infty$ et $\sum \beta_n < \infty$. On peut donc appliquer le théorème de Robbins-Siegmund pour conclure que:
    \[
        \lim_{n \to \infty} z_n = 0 \quad \text{p.s.},
    \]
    c'est-à-dire $\lim_{n \to \infty} \theta_n = \theta^*$ p.s.
    \qed
\end{proof}




% \begin{theorem}[Robbins-Siegmund]
%     Soient $(z_n)_{n \in \N}, ( \alpha_n)_{n \in \N}, ( \beta_n)_{n \in \N}$ et 
%     $ ( \gamma_n)_{n \in \N}$ des processus stochastiques adaptés pour une filtration 
%     $ (\mathcal{F}_n)_{n \in \N}$ tels que : 
%     \begin{enumerate}[label=\roman*)]
%         \item $ \sum_{n \geqslant 1} \alpha_n = \infty $ p.s 
%         \item $ \sum_{n \geqslant 1} \beta_n < \infty $ p.s 
%         \item $ \sum_{n \geqslant 1} \gamma_n < \infty $ p.s 
%     \end{enumerate}
%     et 
%         $$ \mathbb{E}(z_{n+1} \mid \mathcal{F}_n) \leqslant (1 - \alpha_n + \beta_n) z_n + \gamma_n $$ 
%     alors $ \lim_{n \to \infty} z_n = 0$.
% \end{theorem}

% \begin{proposition}
%     Soit $( \alpha_n)_{n \in \N}$ et $ ( \gamma_n)_{n \in \N}$ sont deux suites positives tells 
%     que : 
%         $$ \sum_{n \geqslant 0} = \infty \quad \text{et} \quad \sum_{n \geqslant 0} \gamma_n < \infty $$ 
%     et si $(z_n)_{n \in \N}$ est une suite positive vérifiant : 
%         $$ z_{n+1} \leqslant (1 - \alpha_n)z_n + \gamma_n $$ 
%     alors $ z_n \longrightarrow 0$. 
% \end{proposition}

% \begin{proof}
%     On a pour tout $n \in \N$:
%         $$ z_{n+1} - z_n \leqslant \gamma_n \alpha_n z_n $$ 
%     Par télescopage, on obtient : 
%         \begin{align*}
%             & z_{n} - z_0 \leqslant \sum_{k=0}^{n-1} \gamma_k - \sum_{k=0}^{n-1} \alpha_k z_k \\ 
%             \iff & \sum_{k=1}^{n-1} \alpha_k z_k \leqslant \sum_{k=0}^{n-1} \gamma_k + z_0 - z_n \\
%             & \leqslant \leqslant \sum_{k=0}^{n-1} \gamma_k + z_0 \leqslant \sum_{k=0}^{ \infty} \gamma_k + z_0 
%         \end{align*}
%     Par conséquent, la série $\sum \alpha_k z_k$ converge donc $ z_k \longrightarrow 0$.
% \end{proof}

% \begin{lemma}
%     Sans perdre de généralité, on peut considérer $ \beta_n = 0$ pour 
%     tout $n \in \N$.
% \end{lemma}

% \begin{proof}
%     On a : 
%         \begin{align*}
%             \mathbb{E} \left( \frac{z_{n+1}}{ \prod_{j=1}^{n} (1 + \beta_j)} \mid \mathcal{F}_n\right) &= \frac{ \mathbb{E}(z_{n+1} \mid \mathcal{F}_n)}{ \prod_{j=1}^{n} (1 + \beta_n)} \\ 
%             & \leqslant \frac{(1 + \beta_n - \alpha_n)}{  \prod_{j=1}^{n} (1 + \beta_j) } z_n + \frac{ \gamma_n}{ \prod_{j=1}^{n} (1 + \beta_j)} \\ 
%             &= \left(1 - \frac{ \alpha_n}{ 1 + \beta_n}\right) \frac{z_n}{ \prod_{j=1}^{n} (1 + \beta_j)} - \frac{ \gamma_n}{ \prod_{j=1}^{n} (1 + \beta_j)}
%         \end{align*}
%     Donc : 
%         $$ \mathbb{E}(Y_{n+1} \mid \mathcal{F}_n) \leqslant (1 - \alpha_n') Y_n + \gamma_n' $$ 
%     où $ Y_n = \displaystyle\frac{z_n}{ \prod_{j=1}^{n} (1 + \beta_j)}, \alpha_n' = \displaystyle\frac{ \alpha_n}{1 + \beta_n}$ 
%     et $ \gamma_n' = \displaystyle\frac{ \gamma_n}{ \prod_{j=1}^{n} (1 + \beta_j)}$.
    
%     De plus, $ \sum \alpha_n' = \infty $ p.s car $ \beta_n \longrightarrow 0$ p.s donc $ \alpha_n' \sim \alpha_n$. 
%     D'autre part, $ \sum \gamma_n' < \infty$ p.s car $ \gamma_n' < \gamma_n$.
%     On a alors $ z_n = Y_n  \prod_{j=1}^{n} (1 + \beta_j) $.
%     donc 
%         $$ \lim_{n \to \infty} z_n = \lim_{n \to \infty} Y_n  \underbrace{\prod_{j=1}^{n} (1 + \beta_j) }_{ < \infty}$$
% \end{proof}

% \begin{proof}[Preuve - EXAMEN !!!]
%     On suppose que :
%         $$ \mathbb{E}(z_n \mid \mathcal{F}_n) \leqslant ( 1 - \alpha_n) z_n + \gamma_n $$ 
%     On souhaite construire une sur-margingale positive pour étudier notre problème. 
%     Posons donc : 
%         $$ \forall n \in \N, \quad S_n = Z_n + \sum_{k=0}^{n-1} \alpha_k z_k - \sum_{k=0}^{n-1} \gamma_k $$
%     On a : 
%         \begin{align*}
%             \mathbb{E}(S_{n+1} \mid \mathcal{F}_n) &= \mathbb{E}(z_{n+1} \mid \mathcal{F}_n) 
%                 + \mathbb{E} \left( \sum_{k=0}^{n} \alpha_k z_k \mid \mathcal{F}_n \right)
%                 - \mathbb{E} \left( \sum_{k=0}^{n} \gamma_k \mid \mathcal{F}_n \right) \\ 
%             &= \mathbb{E}(z_{n+1} \mid \mathcal{F}_n) + \sum_{k=0}^{n} \alpha_k z_k - \sum_{k=0}^{n} \gamma_k \\ 
%             & \leqslant z_n - \alpha_n z_n + \gamma_n + \sum_{k=0}^{n} \alpha_k z_k - \sum_{k=0}^{n} \gamma_k \\ 
%             & \leqslant z_n \sum_{k=0}^{n-1} \alpha_k z_k - \sum_{k=0}^{n-1} \gamma_k = S_n
%         \end{align*}
%     C'est donc bien une surmartingale par hypothèse. Posons : 
%         $$ \tau_C = \inf_{n \in \N} \left\{ \sum_{k=1}^{n} \gamma_k > C \right\} $$
%     c'est un temps d'arrêt pour $ (\mathcal{F}_n)_{n \in \N}$. De plus, 
%     $(S_{n \wedge \tau_C})_{n \in \N}$ est une surmatingale et : 
%         $$ S_{n \wedge \tau_C} \geqslant - \sum_{k=0}^{n \wedge \tau_c - 1} \gamma_k \geqslant - \sum_{k=0}^{\tau_c - 1} \gamma_k \geqslant - C $$
%     D'où $(S_{n \wedge \tau_C})_{n \in \N}$ est une surmartingale positive. Donc 
%     $(S_{n \wedge \tau_C} + C)_{n \in \N}$ converge p.s.
    
%     Pour tout $ C > 0$, on a : 
%     $$ \lim_{n \to \infty} S_{n \wedge \tau_C} = W_C \quad \text{p.s} $$
%     Or $\{ \tau_C = \infty\} = \left\{ \sum_{k=1}^{\infty} \gamma_k \leqslant C \right\}$ et 
%     $ \sum_{k=1}^\infty \gamma_k < \infty$ p.s donc, presque sûrement, il existe $C \in \R$
%     tel que $ \tau_C = \infty$ et sur $\{ \tau_c = \infty\}$, on a $\lim_{n \to \infty} S_n = W_C$. 
%         $$ \{ \tau_n = \infty \} \subset \{ \omega \in \Omega \mid \lim S_n \text{ existe} \} $$
%     Donc $ \myP(\lim S_n \text{ existe} ) = 1$.

%     Par conséquent $ \left( S_n + \sum_{k=0}^{n-1} \gamma_k \right)_{n \in \N}$ converge, donc 
%     $ \left(z_n + \sum_{k=0}^{n-1} \alpha_k z_k \right)_{n \in \N}$ converge. Or 
%     $ \sum_{k=0}^{\infty} \alpha_k = \infty $ p.s et $ \sum_{k=0}^{\infty} \alpha_k z_k < \infty$ p.s 
%     donc $ \lim_{n \to \infty} z_n = 0 $ p.s.
% \end{proof}

% \subsection{Conditions suffisantes pour la convergence de l'algorithme de RM} 

% Pour rappel, on a posé : 
%     $$ \theta_{n+1} = \theta_n - \alpha_n h( \theta_n) + \alpha_n \Delta M_n $$ 
% où $ \Delta M_n = h( \theta_n) - H( \theta_n, Y_{n+1})$ et $h( \theta) = \mathbb{E}(H( \theta, Y))$. 
% Sous de bonnes hypothèses sur $H$, on a : 
%     $$ \lim_{n \to \infty} \theta_n = \theta^* \quad \text{p.s} $$  
%     $$ \lim_{n \to \infty} \frac{ \theta_n - \theta^*}{ \sqrt{ \alpha_n}} = \mathcal{N}(0, \sigma^2) \quad \text{en loi} $$ 

% \begin{theorem}[Conditions suffisantes]
%     Si il existe $ \theta^*$ tel que $h( \theta) ( \theta - \theta^*) \geqslant \mu \| \theta - \theta^* \|^2$ 
%     (coercivité). Si il existe $A, B > 0$ tels que $ \| h( \theta_n) \| ^2 \leqslant A + B \| \theta_n \|^2$ et que 
%     $ \mathbb{E}( \| h( \theta) - H( \theta, Y) \|^2) \leqslant K$ (contrôle du bruit).

%     Alors $\lim_{n \to \infty} \theta_n = \theta^*$ p.s.
% \end{theorem}

% \begin{proof}
%     Calculons : 
%         \begin{align*}
%             & \mathbb{E}( \| \varepsilon_{n+1} - \theta^* \|^2 \mid \mathcal{F}_n) \\ 
%             &= \mathbb{E}( \| \theta_n - \alpha_n h( \theta_n) + \alpha_n \Delta M_n - \theta^* \|^2 \mid \mathcal{F}_n) \\ 
%             &= \mathbb{E}( \| \theta_n - \theta^*\|^2 + \alpha_n^2 h( \theta_n)^2 + \alpha_n^2 \| \Delta M_n \|^2 \\
%             & \quad - 2 \alpha_n ( \theta_n - \theta^*) . h( \theta_n) + 2 \alpha_n ( \theta_n - \theta^*) . \nabla M_n - \alpha_n^2 h( \theta_n) . \Delta M_n \mid \mathcal{F}_n) \\ 
%             &= \| \theta_n - \theta^* \|^2 + \alpha_n^2 \| h( \theta_n) \|^2 + \alpha_n^2 \mathbb{E}( \Delta M_n \mid \mathcal{F}_n) - 2 \alpha_n ( \theta_n - \theta^*) . h( \theta_n) 
%         \end{align*}
%     On obtient donc : 
%         \begin{align*}
%             \mathbb{E}( \| \theta_{n+1} - \theta^*\|^2 \mid \mathcal{F}_n) & \leqslant \| \theta_n - \theta^* \|^2 (1 - 2 \alpha_n \mu) + \alpha_n^2 [ K + A + B \| \theta_n - \theta^*\|^2]
%         \end{align*}
%     En posant $ z_n = || \theta_n - \theta^*\|^2$, on a : 
%         \begin{itemize}
%             \item $ " \alpha_n" = 2 \alpha_n \mu$ 
%             \item $ " \beta_n" = \alpha_n^2 \beta_n$ 
%             \item $ " \gamma_n" = \alpha_n^2 (A + K)$ 
%         \end{itemize}
%     Donc 
%         $$ \mathbb{E}(z_{n+1} \mid \mathcal{F}_n) \leqslant (1 - \alpha_n + \beta_n) z_n + \gamma_n $$
%     et on applique le résultat précédent. 
% \end{proof}

% \section{Cas pratique}

% \subsection{Descente de gradient stochastique}

% On souhaite déterminer le point $x \in \R$ qui minimise la fonction $V(x) = \frac{1}{K} \sum_{k=1}^{K} f_k(x)$.

% L'algorithme du gradient nous donne : 
% \begin{itemize}
%     \item On part de $x_0$. 
%     \item On calcule $ \nabla V(x_0)$. 
%     \item On pose $x_1 = x_0 - \alpha_0 \nabla V(x_0)$.
% \end{itemize}

% Algorithme stochastique : 
% \begin{itemize}
%     \item On part de $x_0$ 
%     \item A chaque étape : 
%         \begin{itemize}
%             \item On tire $I \sim \mathcal{U}( \llbracket 1, K \rrbracket)$. 
%             \item On pose $x_{n+1} = x_n - \alpha_n \nabla f_I(x_n) $.
%         \end{itemize}
% \end{itemize}

% C'est un algorithme de type Robbins Monro.

% \subsection{Régression "inline"}

